%
%Nesta seção serão apresentados os recursos tecnológicos usados para o desenvolvimento do projeto.
%
%
%%\subsection{Figma}
%%
%%O Figma foi escolhido como software de elaboração do protótipo, devido à sua alta robustez como ferramenta de design colaborativo. O mesmo permite a inclusão de padrões e templates que são amplamente utilizados na indústria. Além disso, o Figma oferece a capacidade de criar fluxo de telas, possibilitando a realização de protótipos de alta fidelidade, com integração de fluxo de dados \cite{figma}.
%%
%%\subsection{Draw.io BrModelo Bizagi, Miro (ferramentas de diagramação, essas 4 em conjunto foram utilizadas para desenvolver os mais diversos tipos de diagrama no projeto)
%%\subsection{Miro}
%%
%%A metodologia Lean Inception exige um ambiente colaborativo, tendo isso em mente, o Miro oferece templates para o desenvolvimento do Lean Inception, além de exemplos de utilização de artefatos e tutoriais para o uso. Logo, essa ferramenta se revela essencial para facilitar a aplicação dessa metodologia \cite{caroli2017lean}.
%%Draw.io é uma ferramenta amplamente utilizada para a elaboração de diagramas UML (Unified Modeling Language), permitindo a criação de diagramas estruturados como diagramas de classes e de diagrama de pacotes. Ademais, oferece suporte a ambientes colaborativos, facilitando o desenvolvimento em equipe \cite{drawio}.
%%
%%\subsection{PowerBI}
%%Trata-se de um software dedicado à análise de dados, sendo capaz de fornecer dashboards interativas que facilitam a visualizações de dados por meio de gráficos, relatórios e Key Performance Indicators (KPIs) \cite{powerbi}. A ferramenta possibilita o compartilhamento de relatórios em tempo real e irá auxiliar na elaboração de atas de progressos individuais mensais com base nos formulários diários, além de atas de visão geral de dados que poderão ser acessados pelos administradores do aplicativo.
%
%\subsection{Spring Boot}
%
%O Spring é um framework utilizado no ecossistema Java, o mesmo é voltado para aplicações robustas e escaláveis. Ele possui suporte à injeção de dependências, configuração de rotas, conexão com base de dados e diversas outras funcionalidades. Em especial, o Spring Boot é responsável por gerar o projeto com as configurações pré-definidas, evitando a adição manual das dependências. O mesmo segue o princípio "convention over configuration", que gera as dependências de forma automática \cite{weissmann2014vire}.
%
%\subsection{PostgreSQL}
%
%Como definido na arquitetura do sistema, é necessário um banco relacional para o desenvolvimento deste trabalho. Para tal finalidade, optou-se pela utilização do PostgreSQL, um Sistema Gerenciador de Banco de Dados Relacional (SGBDR) de código aberto amplamente reconhecido pela sua robustez, confiabilidade e escalabilidade.
%
%A escolha do PostgreSQL justifica-se pelas suas diversas funcionalidades avançadas, como suporte à programação de eventos temporais por meio da ferramenta pgAgent, para backups, além de possibilidade de extensão, pois o mesmo possui diversas bibliotecas externas. O PostgreSQL também oferece suporte a \textit{triggers} e \textit{procedures} implementadas em diversas linguagens, como PL/pgSQL, PL/Python, PL/Java, entre outros, o que proporciona maior flexibilidade e customização caso seja necessário \cite{postgresql_2025}.
%
%\subsection{Git e GitHub}
%
%A fim de gerenciar o versionamento do projeto, é necessária uma ferramenta capaz de realizar essa tarefa, para isso se dá a utilização do Git, o mesmo é capaz de rastrear alterações, facilitar o trabalho em equipe e manter o histórico de desenvolvimento. Além disso, também é indispensável uma plataforma de hospedagem de repositórios, para isso o Github se adequa, sendo ele a maior e mais popular plataforma de hospedagem de repositórios \cite{githubdocs}.



Nesta seção serão apresentados os recursos tecnológicos usados para o desenvolvimento do projeto, apresentados na Tabela \ref{tab:tecnologias}.

\begin{center}
    \begin{longtable}{|>{\raggedright\arraybackslash}p{3.5cm}|p{11.5cm}|}
        \caption{Recursos Tecnológicos} \label{tab:tecnologias} \\
        \hline
        \textbf{Tecnologia} & \textbf{Uso no projeto} \\
        \hline
        \endfirsthead

        \hline
        \textbf{Tecnologia} & \textbf{Uso no projeto} \\
        \hline
        \endhead

        Figma e Figma Make & Elaboração do protótipo de alta fidelidade e geração prevista do software web a partir do design \cite{figma}. \\
        \hline
        Ferramentas de diagramação (Miro, Draw.io, BrModelo e Bizagi) & Elaboração dos diversos diagramas do projeto, como Lean Inception \cite{caroli2017lean}, UML \cite{drawio}, entidade-relacionamento e processos de negócio (BPMN). \\
        \hline
        Spring (Spring Boot, Spring Data e Hibernate) & Framework Java utilizado no back-end, incluindo configuração automática, ORM e demais APIs do ecossistema \cite{weissmann2014vire}. \\
        \hline
        Git e GitHub & Versionamento e hospedagem dos repositórios do projeto, sendo eles Monografia, Front e Service \cite{githubdocs}. \\
        \hline
        Claude Code & Apoio ao desenvolvimento do software, auxiliando na escrita, revisão e depuração de código. \\
        \hline

        \multicolumn{2}{c}{\small Fonte: autoria própria.} \\
    \end{longtable}
\end{center}